\appendix

\section{COPASI-Modell}

Im Folgenden werden die für die Implementierung relevanten
COPASI-Definitionen tabellarisch zusammengefasst. Das vollständige Projekt
inklusive COPASI-Modell und Auswertungsdateien ist öffentlich auf GitHub
verfügbar: \url{https://github.com/PatrickStahl/MolekulareAlgorithmen}

\subsection{Benutzerdefinierte Funktionen}

\begin{table}[h]
    \centering
    \begin{tabular}{c|c|c}
        Name & Typ & Mathematische Beschreibung \\
        \hline
        \texttt{v\_square} & userdefined & \texttt{S*S} \\
        \texttt{v\_scaled\_Q} & userdefined & \texttt{c * Q} \\
        \texttt{v\_const\_theta} & userdefined & \texttt{theta\_ref} \\
    \end{tabular}
    \caption{Benutzerdefinierte kinetische Funktionen in COPASI.}
    \label{tab:copasi_functions}
\end{table}

\subsection{Spezies}

\begin{table}[h]
    \centering
    \begin{tabular}{c|c|c|c}
        Name & Kompartiment & Typ & Anfangskonzentration \\
        \hline
        \texttt{S} & \texttt{cell} & fixed & \(s\), exemplarisch \(2\) \\
        \texttt{Q} & \texttt{cell} & reactions & \(0\) \\
        \texttt{F} & \texttt{cell} & reactions & \(0\) \\
        \texttt{Theta\_n} & \texttt{cell} & reactions & \(0\) \\
    \end{tabular}
    \caption{Im Modell verwendete Spezies.}
    \label{tab:copasi_species}
\end{table}

\subsection{Globale Größen}

\begin{table}[h]
    \centering
    \begin{tabular}{c|c|c}
        Name & Typ & Wert \\
        \hline
        \texttt{c\_area} & fixed & \(0.716446108\) \\
        \texttt{theta\_ref} & fixed & \(0.819545628\) \\
    \end{tabular}
    \caption{Globale Modellparameter.}
    \label{tab:copasi_global_quantities}
\end{table}

\subsection{Reaktionen}

\begin{table}[H]
    \centering
    \begin{tabular}{c|c|c}
        Name & Reaktionsschema & Rate Law \\
        \hline
        \texttt{R\_Q\_prod} & \(\varnothing \rightarrow Q;\ S\) & \texttt{v\_square} \\
        \texttt{R\_Q\_deg} & \(Q \rightarrow \varnothing\) & Mass action, \(k=1\) \\
        \texttt{R\_F\_prod\_1} & \(\varnothing \rightarrow F;\ Q\) & \texttt{v\_scaled\_Q} \\
        \texttt{R\_F\_prod\_c} & \(\varnothing \rightarrow F;\ Q\) & \texttt{v\_scaled\_Q} \\
        \texttt{R\_F\_deg} & \(F \rightarrow \varnothing\) & Mass action, \(k=1\) \\
        \texttt{R\_Theta\_prod} & \(\varnothing \rightarrow \Theta_n\) & \texttt{v\_const\_theta} \\
        \texttt{R\_Theta\_deg} & \(\Theta_n \rightarrow \varnothing\) & Mass action, \(k=1\) \\
    \end{tabular}
    \caption{Reaktionen des COPASI-Modells. Spezies hinter dem Semikolon
    wirken als Modifier und werden durch die jeweilige Reaktion nicht
    verbraucht.}
    \label{tab:copasi_reactions}
\end{table}

\subsection{Parameterzuordnungen der Reaktionen}

Da die Funktion \texttt{v\_scaled\_Q} mehrfach verwendet wird, sind die
jeweiligen Parameterzuordnungen in Tabelle
\ref{tab:copasi_parameter_mapping} separat aufgeführt.

\begin{table}[h]
    \centering
    \begin{tabular}{c|c|c}
        Reaktion & Funktionssymbol & Zugeordnetes Modellobjekt / Wert \\
        \hline
        \texttt{R\_Q\_prod} & \texttt{S} & Spezies \texttt{S} \\
        \texttt{R\_F\_prod\_1} & \texttt{Q} & Spezies \texttt{Q} \\
        \texttt{R\_F\_prod\_1} & \texttt{c} & \(1\) \\
        \texttt{R\_F\_prod\_c} & \texttt{Q} & Spezies \texttt{Q} \\
        \texttt{R\_F\_prod\_c} & \texttt{c} & globale Größe \texttt{c\_area} \\
        \texttt{R\_Theta\_prod} & \texttt{theta\_ref} & globale Größe \texttt{theta\_ref} \\
    \end{tabular}
    \caption{Relevante Parameterzuordnungen der benutzerdefinierten
    kinetischen Funktionen.}
    \label{tab:copasi_parameter_mapping}
\end{table}

Die Zuordnung in Tabelle \ref{tab:copasi_parameter_mapping} stellt
insbesondere sicher, dass für die Flächenausgabe zwei parallele
Produktionsreaktionen mit den Raten \([Q]\) und \(c[Q]\) verwendet werden.
Damit ergibt sich insgesamt

\[
    \frac{d[F]}{dt}
    =
    [Q]+c[Q]-[F].
\]

Alle verwendeten Ratenkonstanten besitzen Werte kleiner oder gleich \(1\).
